% ===== §8 The Justice Theory of Choice =====
\section{The Justice Theory of Choice}
\label{p16:sec:justice}

\noindent
The theory of relational luxury developed across the preceding sections is not merely descriptive; it carries normative implications that constitute a distinctive justice theory. The choice of one form of luxury over another---the choice of the flower over the diamond, or the diamond over the flower, or one mode of giving over another---is not merely a matter of personal taste but a value declaration with justice dimensions. This section develops the justice theory of choice in three movements: the problem of subject-embedding misrecognition, the problem of equivalent agency, and the constructive account of selection justice.

The justice theory developed here is, by the standards of the dominant frameworks, unusual: it does not concern the distribution of goods, the protection of rights, or the maximisation of welfare, but the integrity of existence-attestation---the question of whether what is presented as carrying a subject's existence genuinely carries it, and whether the substitution of one form of existence-attestation for another is just. This is a justice of the real register: a justice concerned with the genuineness of existence-attestation rather than the fairness of distribution.

\subsection{Subject-Embedding Misrecognition: The True Problem Domain}
\label{p16:subsec:misrecognition}

\noindent
The justice theory of choice begins not with the question of what to choose but with a prior epistemological question: \emb{is the subject-embedding that an object presents itself as carrying genuinely present, and is it accurately recognised?} This question identifies the true problem domain of the justice theory: the domain of subject-embedding misrecognition, in which the relationship between an object's apparent and actual coupling degree is distorted.

Subject-embedding misrecognition takes three forms, corresponding to three loci of distortion.

\emb{The giver's self-misrecognition}: The first form occurs when the giver misrecognises their own existence-attestation---when they believe that they have embedded their existence in an object that, in fact, carries low coupling degree. The paradigmatic case is the giver who believes that economic expenditure is equivalent to existential embedding: who delegates the purchase of an expensive gift to an assistant and sincerely believes that the resulting gift expresses the depth of their care, because they have confused the economic cost of the gift with the existential cost of genuine attentiveness.

This self-misrecognition is the most insidious form because it operates beneath the giver's awareness. The giver is not consciously simulating care; they genuinely believe that they are expressing it. The consumer ideology that systematically equates economic value with relational value---that trains its participants to read the price of a gift as a measure of the love it expresses---makes this self-misrecognition natural and pervasive. The giver who has internalised this ideology cannot easily distinguish their economic expenditure from their existential investment, because the ideology has collapsed the distinction.

\emb{The receiver's induced misrecognition}: The second form occurs when the receiver is induced to misrecognise the coupling degree of an object---when the sign exchange value of an expensive object induces the receiver to believe that it carries high coupling degree when, in fact, it does not. This is the precise analogue, in the domain of luxury, of the forged trust analysed in \pinkref{Paper~XIII} \parencite{paperxiii}: the inducement of the receiver's generative model to lower its prediction error---to trust, to feel cared for, to believe in the giver's existential investment---without the provision of a genuine high-fidelity signal of that investment.

The mechanism of induced misrecognition is the sign exchange value system itself: because consumer culture has established the convention that expensive gifts express serious care, the receiver who has internalised this convention reads the expensive gift as evidence of care, even when the gift carries low coupling degree. The expensive diamond, purchased without genuine attentiveness, induces the belief in existential investment through its sign exchange value alone---``he spent so much, he must really love me''---and this induced belief is a misrecognition of the actual coupling degree.

\emb{The social collective misrecognition}: The third form is structural rather than individual: it is the collective misrecognition, embedded in the social symbolic order, that equates economic value with coupling degree as a general matter. This collective misrecognition is the cultural condition that makes the individual forms possible: it is because the society as a whole has come to read economic value as a proxy for relational value that individual givers can misrecognise their own attestation and individual receivers can be induced to misrecognise what they receive.

The social collective misrecognition has a specific and damaging consequence: it renders the flower invisible. Within a symbolic order that equates economic value with coupling degree, the relationally luxurious but economically modest gift---the flower, the handmade object, the found thing carrying a shared memory---cannot be recognised as luxury at all. It is systematically devalued, dismissed as charming but insignificant, because the collective misrecognition has no category for high coupling degree at low economic value. The social collective misrecognition is thus the cultural mechanism by which the third doctrine's central insight---that coupling degree is the true measure of intimate luxury---is rendered culturally invisible.

\begin{claim}[The three forms of subject-embedding misrecognition]
Subject-embedding misrecognition---the distortion of the relationship between an object's apparent and actual coupling degree---takes three forms: the giver's self-misrecognition (confusing economic expenditure with existential investment), the receiver's induced misrecognition (being led by sign exchange value to perceive coupling degree that is not present), and the social collective misrecognition (the cultural equation of economic value with coupling degree that renders the flower invisible). These three forms are mutually reinforcing: the social collective misrecognition is the cultural condition that makes the individual forms possible, and the individual forms reproduce the collective misrecognition.
\end{claim}

\subsection{The Problem of Equivalent Agency}
\label{p16:subsec:equivalent_agency}

\noindent
The second movement of the justice theory addresses a problem that is more subtle than misrecognition and that constitutes the theory's most original contribution: the problem of \emb{equivalent agency}.

The problem arises as follows. A subject wishes to provide an existence-attestation $A$---the immediate, present, full-attention embedding of the flower found on the path, say. But the subject cannot provide $A$: they are not in the place where the flower grows, they do not have the time for the walk on which it would be found, the spatiotemporal conditions for $A$ are not available to them. Instead, the subject provides $B$---a different form of existence-attestation, such as the long-saved ring, which embeds a different temporal structure (sustained investment over time rather than immediate presence). The subject substitutes $B$ for $A$: they use one form of existence-attestation in place of another that they cannot provide.

\emb{Is this substitution just?}

The question is genuinely difficult, because $A$ and $B$ are not equivalent in their structure. They embed different temporal structures, produce different qualities of coupling, and carry different existential meanings. From the standpoint of \emb{structural justice}---a justice concerned with the structural identity of what is provided---the substitution of $B$ for $A$ is a form of alienation: $B$ is not $A$, and the substitution presents $B$ as if it were equivalent to $A$ when it is structurally heterogeneous. Certain justice frameworks, attentive to this structural heterogeneity, will judge equivalent agency as unjust: as the passing-off of one thing for another, the substitution of a different (and possibly lesser) existence-attestation for the one that was due.

But from the standpoint of \emb{generative justice}---the framework developed in \pinkref{Paper~XV} \parencite{paperxv}, concerned with whether the generative process is just rather than whether the products are structurally identical---the substitution may be just, under specific conditions. The generative justice analysis asks not ``is $B$ structurally identical to $A$?'' (it is not) but ``is the generative process by which $B$ was provided in place of $A$ a just process?'' And this question can have an affirmative answer even when $A$ and $B$ are structurally heterogeneous.

\subsubsection{The Four Necessary Conditions}
\label{p16:subsubsec:four_conditions}

\noindent
Equivalent agency is just in the generative sense only if four conditions are satisfied. These conditions are \emb{necessary but not sufficient}---a crucial qualification developed below.

\emb{Condition 1: Symmetric generability}. In ideal conditions---conditions without the real constraints that force the substitution---both $A$ and $B$ would be within the subject's generative range. The subject is not substituting $B$ for $A$ because they are incapable of $A$ in principle, or because they do not understand what $A$ would require, but because specific real circumstances prevent $A$ in this instance. The symmetry condition ensures that the substitution is not a permanent incapacity masquerading as a temporary constraint: the subject who could never provide $A$---who has no capacity for the immediate, present, full-attention embedding that $A$ requires---is not engaged in just equivalent agency when they provide $B$; they are providing the only thing they can provide, which is a different matter.

\emb{Condition 2: Real constraint}. The circumstances that force the substitution of $B$ for $A$ are genuine, not pretextual. The subject genuinely cannot provide $A$---because of real limitations of time, space, or resources---rather than choosing $B$ over $A$ out of laziness, indifference, or a preference for the social legibility of $B$. The real constraint condition distinguishes genuine equivalent agency from the use of constraint as an excuse: ``I was too busy to find the flower, so I bought the ring'' is just equivalent agency only if the busyness is a genuine constraint rather than a rationalisation of the preference for the more socially legible gift.

\emb{Condition 3: Genuine embedding in $B$}. The substitute $B$ genuinely carries the subject's existence-attestation; it is not a mere economic substitute but a real, if different, form of subject-embedding. The long-saved ring is a just substitute for the flower only if it genuinely embeds the subject's existence---the years of saving, the sustained anticipatory attention, the specific choice for this specific person---rather than merely carrying the sign exchange value that the consumer convention associates with serious commitment. If $B$ is a genuine existence-attestation (even of a different structure from $A$), the substitution can be just; if $B$ is a mere economic substitute with low coupling degree, the substitution is not equivalent agency but the replacement of existence-attestation with sign exchange value.

\emb{Condition 4: Transparency}. The receiver can recognise that what they are receiving is $B$ rather than $A$---that the existence-attestation they are receiving has the structure of sustained investment rather than immediate presence, and that it is offered in place of the $A$ that the circumstances prevented. Transparency does not require explicit verbal disclosure (``I would have brought you a flower but I was too busy, so here is a ring''), which would undermine the relational quality of the gift; it requires that the gift be offered and received in a way that does not induce the misrecognition of $B$ as $A$ or of $B$ as a higher-coupling gift than it is. Transparency is the condition that connects equivalent agency to the analysis of misrecognition: just equivalent agency does not involve inducing the receiver to misrecognise the coupling degree or structure of what they receive.

\subsubsection{Necessary But Not Sufficient}
\label{p16:subsubsec:not_sufficient}

\noindent
The four conditions are necessary for just equivalent agency: their violation entails injustice. But they are \emb{not sufficient}: their satisfaction does not guarantee full justice. This qualification is essential and was, in the development of this theory, a correction of an initial over-simplification. Even when all four conditions are met, residual justice questions remain.

\emb{The residual structural heterogeneity}: Even when the four conditions are satisfied, $A$ and $B$ remain structurally heterogeneous: $B$ produces a different quality of coupling from $A$, and this difference is real. Just equivalent agency does not dissolve this heterogeneity; it renders the substitution permissible despite the heterogeneity. The structural difference remains as a \emb{residual justice debt}---not a debt that must be repaid, but a difference that must be acknowledged rather than dissolved. The subject who provides $B$ in place of $A$, even justly, owes the acknowledgement that $B$ is not $A$---that something of the immediate presence that $A$ would have embedded is genuinely absent from $B$, however genuine $B$'s own embedding may be.

\emb{The graded nature of the conditions}: The four conditions are not binary (satisfied or not satisfied) but graded (satisfied to a degree). The real constraint can be more or less genuine; the embedding in $B$ can be more or less deep; the transparency can be more or less complete; the symmetric generability can be more or less robust. Even when all four conditions are ``satisfied'' in the sense of being present, their degree of satisfaction affects the justice of the substitution. A substitution in which all four conditions are robustly satisfied is more just than one in which they are barely satisfied, even though both technically meet the necessary conditions.

\emb{The receiver's free acceptance}: Just equivalent agency requires not only the four conditions on the giver's side but the receiver's free acceptance of $B$ in place of $A$. The receiver who would have genuinely preferred $A$, and who accepts $B$ only because they feel they have no choice---because refusing would damage the relation, because the social convention demands acceptance, because the giver's evident investment in $B$ makes refusal feel ungrateful---has not freely accepted the substitution. The justice of equivalent agency therefore depends on conditions beyond the giver's control: it depends on the receiver's genuine, uncoerced acceptance of the substituted form.

\emb{The systemic effects of habitual substitution}: Even a single just instance of equivalent agency may contribute, if it becomes part of a systematic pattern, to a relational dynamic that is unjust in aggregate. If the substitution of $B$ for $A$ becomes habitual---if the subject always provides the sustained-investment form and never the immediate-presence form, always the ring and never the flower---then even if each individual substitution meets the four conditions, the cumulative effect may be the systematic absence of the immediate-presence embedding that $A$ provides. The justice of equivalent agency must therefore be assessed not only at the level of the individual act but at the level of the relational field's overall dynamics.

\begin{claim}[Equivalent agency and its justice conditions]
The substitution of one form of existence-attestation $B$ for another $A$ that the subject cannot provide---equivalent agency---is structurally heterogeneous (\,$A$ and $B$ produce different qualities of coupling, and structural justice frameworks will register this as alienation\,), but may be just in the generative sense if four necessary conditions are met: symmetric generability, real constraint, genuine embedding in $B$, and transparency. These conditions are necessary but not sufficient: even when all four are satisfied, residual justice questions remain---the acknowledgement of structural heterogeneity, the graded degree of the conditions' satisfaction, the receiver's free acceptance, and the systemic effects of habitual substitution. Justice in equivalent agency is therefore not a threshold that is crossed once and for all but an ongoing relational achievement.
\end{claim}

\subsection{Selection Justice: Choice as Value Declaration}
\label{p16:subsec:selection_justice}

\noindent
The third movement of the justice theory addresses the constructive question: given the analysis of misrecognition and equivalent agency, what does it mean to choose justly in the domain of relational luxury?

The foundational insight is that \emb{choice in the domain of luxury is a value declaration}. Above the threshold of material necessity, every choice of how to express intimate care through objects reveals the chooser's understanding of value: what they take to be real wealth, what they take to be worth investing in, what form of luxury they take to be genuine. The choice between the flower and the diamond, or between genuine and proxy attentiveness, is not a neutral matter of taste but a declaration of what the chooser believes intimate luxury to be.

\emb{The choice oriented by sign exchange value}: The choice that takes sign exchange value as its primary criterion---that chooses the gift for its social legibility, its position in the hierarchy of economic value, its capacity to signal commitment to the social world---is a value declaration that accepts the symbolic order of consumer society. It anchors relational value in the exchangeable symbolic register, using the language of social difference to express the depth of the intimate bond. This is not simply an error; it is a choice with a cost: the cost of drawing the irreducibly particular relational value into the exchangeable symbolic system, and thereby introducing the commodification mechanism analysed in \pinkref{Paper~XV} \parencite{paperxv}.

\emb{The choice oriented by coupling degree}: The choice that takes coupling degree as its primary criterion---that chooses the gift for its capacity to carry genuine existence-attestation, with attention to the specific person and the specific moment rather than to the social legibility of the gift---is a value declaration of a different kind. It is a \emb{justice stance}: a refusal to use exchangeable symbols to bear the inechangeable relation, an insistence on the irreducibility of relational value, a reservation of space for the real register within the symbolic logic of consumer society.

We call this a justice stance, rather than merely an ethical preference, because it concerns the integrity of existence-attestation against the systematic pressure of the social collective misrecognition. To choose coupling degree over sign exchange value, in a cultural environment that systematically equates economic value with relational value, is to resist the collective misrecognition that renders the flower invisible---to insist, against the pressure of the symbolic order, that the true measure of intimate luxury is the genuine embedding of one subject's existence in the life of another.

\subsection{The Dialectic of Selection: No Pure Relational Luxury}
\label{p16:subsec:no_pure}

\noindent
The justice theory of selection must, however, acknowledge a dialectical complication: there is no pure relational luxury, no object that is entirely free of the sign exchange value dimension. Every object given in an intimate relation carries, alongside its coupling degree, some position in the symbolic order---some social legibility, some sign exchange value, some location in the hierarchy of economic value. Even the flower carries a sign: it signifies, within the symbolic order, a certain kind of romantic gesture, a certain relationship to nature and to convention.

This means that just selection is not the selection of a pure relational luxury object---there is no such object---but the selection that takes coupling degree as primary while acknowledging the unavoidable presence of the sign exchange value dimension. The just chooser does not pretend that their gift is free of social meaning; they ensure that the social meaning is subordinate to the relational meaning, that the coupling degree is the primary motivation and the sign exchange value a secondary (and acknowledged) consequence.

The dialectic of selection therefore takes the form not of a choice between pure relational luxury and pure sign exchange value---a choice that the actual structure of objects does not permit---but of a choice about \emb{primacy}: which dimension is the primary motivation, which dimension governs the act of giving, which dimension the giver attends to and cultivates. Just selection is the selection in which coupling degree is primary, sign exchange value is secondary and acknowledged, and the existence-attestation embedded in the gift is genuine rather than simulated.

\subsection{The Reconstruction of Gift Ethics}
\label{p16:subsec:gift_ethics}

\noindent
The justice theory of choice issues in a reconstructed ethics of the gift in intimate relations. The reconstructed gift ethics has several components, each following from the analysis of the preceding sections.

\emb{The primacy of genuine embedding}: The just gift is the gift that carries genuine existence-attestation---in which the giver's spatiotemporal structure is genuinely embedded, whether through the immediate presence of $A$ or the sustained investment of $B$. The just gift cannot be delegated in its essential dimension: while the mechanics of acquisition can sometimes be delegated, the existence-attestation cannot, because it is constituted by the giver's own spatiotemporal investment. The proxy-purchased gift fails the just gift criterion not because delegation is always wrong but because the delegation of the essential existence-attestation empties the gift of its coupling degree.

\emb{The intentional orientation to the specific other}: The just gift is oriented to the specific particularity of this specific relational other---chosen because of what this person is, with attention to their specific character, history, and preferences, rather than because of the gift's position in the hierarchy of economic value. The intentional intensity dimension of coupling degree is the ethical heart of the just gift: the gift that is chosen for this person, rather than chosen as the most impressive option, is the gift that carries genuine intentional embedding.

\emb{The transparency of the attestation's structure}: The just gift is given in a way that does not induce the misrecognition of its coupling degree or structure---that does not present $B$ as $A$, or low coupling degree as high, or sign exchange value as existence-attestation. The transparency requirement is the ethical expression of the analysis of misrecognition: the just giver does not exploit the social collective misrecognition to induce the receiver to perceive coupling degree that is not present.

\emb{The non-possessive mode of giving}: The just gift is given in the non-possessive mode of \emph{xuande}---given as an offering to the relational field rather than as a claim on the receiver, as an expression of care rather than as a mechanism of indebtedness or a display of power. The just gift does not function as a phallic signifier that positions the giver as the one who has and the receiver as the one who lacks; it functions as a contribution to the relational field's GRW, offered without the expectation of equivalent return.

\emb{The poetic mode of receiving}: The just reception---the complement of the just gift---is the reception that receives the gift poetically rather than through settlement: that opens the gift's meaning rather than closing it, that recognises the existence-attestation it carries rather than reading only its sign exchange value, that completes the coupling through genuine recognition rather than through the conventional acknowledgment that the social order prescribes.

\begin{claim}[The reconstructed gift ethics]
The just gift in intimate relations carries genuine existence-attestation (the primacy of genuine embedding), is oriented to the specific particularity of the relational other (intentional intensity), is given transparently with respect to its coupling degree and structure (no induced misrecognition), is offered in the non-possessive mode of \emph{xuande} (not as phallic signifier or mechanism of indebtedness), and is received poetically rather than through settlement (the completion of coupling through genuine recognition). The justice of the gift is not a property of the object's economic value but of the integrity of the existence-attestation it carries and the genuineness of the recognition with which it is received.
\end{claim}
